\section{Robustness Checks}\label{sec:robustness}

This section collects the methodological checks referenced throughout the paper. Two of them changed reported results. The recalibrated cross-design variant of Section~\ref{sec:robustness-crossdesign} exceeds my ex-ante threshold and qualifies the household-choice reading at the aggregate level. The off-window floor re-measurement of Section~\ref{sec:robustness-floor} is where the paper's headline marginal and level are actually derived ($+9.2 \to +5.6$ points; 97.9\% $\to$ 91.3\% shared). It demotes the in-window calibration to an in-sample point. On the shared basis the paper quotes throughout, the hazard framework undershoots the benchmark at the headline calibration (91.3\%); it overshoots only on the standalone basis (107.0--121.5\%). The remaining checks confirm the central decomposition.

Four of this section's checks are mechanics rather than results, and the online appendix collects them: the behavioral-extension mechanics (Appendix~\ref{sec:robustness-extensions}), the weighted-average-life quantification (Appendix~\ref{sec:robustness-wal}), the fifteen-year fold-in and the rate-gap-units correction (Appendix~\ref{sec:robustness-pipeline}), and the seasonal-floor program (Appendix~\ref{sec:robustness-seasonalfloor}). The calibration box already carries the seasonal-floor verdict. What a seasonally varying floor moves is floor dispersion under the hard maximum, not cyclicality. Monthly timing earns no credential there. The frozen timing rule is cleared both by a wrong-month rotation and by the $\beta_1 = 0$ null itself.

\subsection{Scale Sensitivity}\label{sec:robustness-scale}

The ABM's 10,000-household population and the hazard framework's 75,000-loan sample differ in both scale and data source. To test whether population size alone drives the gap, I re-ran the ABM's core household mechanic at $N = 10{,}000$ and $75{,}000$, 50 seeds each. That mechanic is the production code's exponential, dollar-scale mobility-desire heterogeneity mapped through a rate-gap-driven CPR function. Mean simulated CPR is 15.616\% at both, a difference of 0.0003 percentage points. The 95\% confidence interval narrows 2.4$\times$ at the larger $N$. That is about 12\% short of the 2.74$\times$ an idealized $\sqrt{N}$ law predicts ($\sqrt{7.5} \approx 2.74$), consistent with a non-i.i.d.\ component of seed noise; I do not lean on the exact ratio. Population size changes that mechanic's precision, not its central tendency. Two limits, both large, so I state the conclusion conditionally. The check ran the household cost-benefit mechanic \emph{in isolation}. Its 15.616\% sits 3.9 points (about a third) above the production ABM's 11.68\%, so the object shown scale-invariant is not the estimator whose recovery is at issue. And no committed script, artifact or gate stood behind it; the figure survived from an earlier draft. A pre-committed rerun on the full production behavioral specification addresses both (run \texttt{production\_scale\_test}, fifty seeds per size). It carries the 43\% front-end DTI wall, the discrete twenty-percent wait-and-see freeze trait, income-scaled curtailment, eleven-cohort multi-vintage weighting with its fifteen- and thirty-year term split, and a settlement-lag kernel --- none of them present in the isolated mechanic. It recovers 12.66\% of the benchmark at $N = 10{,}000$ against 12.18\% at $N = 75{,}000$: $0.48$ points against a fifty-seed standard deviation of $3.25$ points, well inside one standard deviation. Mean simulated CPR is 11.60\% and 11.64\%, close to the native-gate run's 11.76\%, not the isolated mechanic's 15.6\%. One dimension of the production specification this run does not carry, and it is the one that reconciles it with Table~\ref{tab:estimators}. The rerun evaluates the fifteen-year book on the native behavioral gate; production folds that book in on structural terms only (Appendix~\ref{sec:robustness-pipeline}). Table~\ref{tab:estimators}'s 13.56\% fifty-seed mean is the production fold-in, this arm's 12.66\% the native gate. The 0.90 points between them are that documented level offset, which the run was specified to expose, and not an effect of population size. Both arms here sit on the same gate, so the scale comparison itself is like-for-like. Seed precision improves with population: the standard deviation falls from $3.25$ to $1.05$ points, a ratio of 3.1 against the same $\sqrt{7.5} \approx 2.74$, and again I do not lean on the exact ratio. Population size therefore does not drive the gap between this estimator and the hazard framework, and that now rests on the production specification, not a fragment of it. It does not follow that the gap is modeling paradigm alone. Resampling cannot close that distinction, and Section~\ref{sec:robustness-crossdesign} takes it up with the completed cross-design pair.

\subsection{Benchmark Sensitivity to Data Source and Cap-Schedule Assumptions}\label{sec:robustness-benchmark}

The \$764.7 billion empirical benchmark (Section~\ref{sec:method-benchmark}) is constructed from the New York Fed's SOMA current-face-value series under the actual phased redemption cap schedule, netting realized roll-off against the cap within the active QT window. I treat net accumulation, not one-sided accumulation of shortfall months only, as the correct construction because it is the only version comparable to the Danish counterfactual of Section~\ref{sec:abm-danish}, where simulated roll-off routinely exceeds the cap; a one-sided construction would overstate the empirical benchmark relative to that counterfactual and bias the recovered percentages reported for both estimators. Researchers extending this work with alternative SOMA vintages or a different cap-schedule assumption should expect the benchmark level, and therefore the recovered percentages in Sections~\ref{sec:abm} and~\ref{sec:hazard}, to shift. The qualitative ranking of the three estimators reported in Table~\ref{tab:estimators} is not sensitive to this choice, since the ABM's shortfall against any plausible benchmark construction remains an order of magnitude smaller than either hazard path's. Section~\ref{sec:method-benchmark} separately addresses the conceptual status of the cap-relative construction itself. One such alternative is now measured. The benchmark nets cash arriving in month $t$ against the cap in force in month $t$, which is what the cap constrains: the Committee set a ceiling on principal not reinvested in a given month, and that is a statement about cash. Agency remittance, however, means the cash arriving in month $t$ was generated by prepayment decisions in months $t$, $t-1$ and $t-2$, on the settlement kernel this paper already applies to simulated roll-off (0.10, 0.60, 0.30 at lags zero through two). Agency remittance is itself a deterministic lag of roughly one month, so that profile is not a description of the remittance cycle: it is an empirical smearing kernel over the months in which a given month's prepayment decisions arrive as cash, and the alternate kernels below price the choice. Convolving the cap schedule forward with that same kernel and re-netting the unchanged realized series gives a settlement-aligned benchmark of \$722.7 billion against the calendar \$764.7 billion, a shift of $-5.5$\% (run \texttt{settlement\_months\_benchmark}). The shift is mechanical, and its size is exactly accounted for. Because the QT window is half-open, cap convolved past its final month is dropped: October 2025 loses its two-month tap and November 2025 its one- and two-month taps, \$42.0 billion of the \$1{,}417.5 billion schedule. The realized series is untouched and both legs sum over the same months, so the benchmark moves by exactly that \$42.0 billion. Two alternate kernels bracketing that profile put the shift at $-\$28.0$ and $-\$52.5$ billion. Aligning the realized leg as well moves the benchmark less, not more. The pre-QT cap is zero, so the cap leg loses mass only at the window's end, whereas realized roll-off is not zero before the window and the two months the kernel reaches back into settle into its start: the realized leg gains \$6.8 billion there against \$23.4 billion leaving at the end, giving back \$16.6 billion of the cap-side \$42.0 billion. Convolving both legs therefore lands the benchmark at \$739.4 billion, a shift of $-\$25.4$ billion or $-3.3$\%, which the same two alternate kernels bracket at $-\$19.0$ and $-\$29.2$ billion (run \texttt{realized\_boundary\_allocation}). That the two edge effects would partially offset was fixed before the run; the sign of the residual deliberately was not, and at 3.3\% of the committed benchmark the boundary treatment stays a disclosed sensitivity rather than a re-basing: the production convention remains calendar on both legs. I retain the calendar convention, for the reason Section~\ref{sec:method-benchmark} gives for retaining the cap at all: it is the only policy-committed quantity available, and it is committed in cash-arrival months. What the cap-only variant establishes is that re-basing moves no dollar quantity in this paper (the identified marginal is \$42.6 billion on either denominator) and moves the percentages only by the common factor $1.058$: the central Path~B level reads 96.6\%, not 91.3\%, the $\beta_1 = 0$ null 90.7\%, not 85.7\%, and the marginal $+5.9$ points, not $+5.6$. The decomposition is unchanged in substance, and the null still recovers the large majority of the benchmark under either convention. The cap-only variant also bears on the path-level timing question every estimator here leaves open. By a cash-arrival object I mean reported roll-off smeared over the months in which a month's prepayment decisions arrive as cash and, in its CPR form, clipped by a back-out whose non-negativity rule pins four months to exactly zero and posts the displaced mass into the following month's reading (Section~\ref{sec:method-benchmark}). If the monthly series each estimator's path is scored against is such an object, then a monthly path is not a behavioral object, and the timing question may be ill-posed. I state that as a candidate account, not a finding: reporting mechanics and behavior arrive in the same series, and nothing here separates them. A related window-construction pitfall merits a reproducibility note: accumulating from the QT start date without an upper bound admits eight post-QT months (December 2025 onward). Where the dollar aggregation already applies the bounded active-QT mask, as it does throughout this paper, only the path statistics drift (the headline CPR means by roughly 0.3--0.5 percentage points); where it does not, the benchmark level itself moves. That asymmetry (path statistics silently sensitive to a window filter the masked dollar aggregates are robust to) is why every headline figure in this paper cites a single tagged run manifest (footnote~\ref{fn:manifest}).
